\documentclass[12pt]{article}
\usepackage[utf8]{inputenc}
\usepackage[T1]{fontenc}
\usepackage{hyperref}
\title{Tjupany Language Resources for K-12 Curriculum Development}
\author{Compiled by Claude Sonnet 5 (claude-sonnet-5) for Andr\'as Kornai}
\date{}
\begin{document}
\maketitle

\section{Introduction}
Tjupany (also spelled Tjupan; ISO 639-3 tjp, Glottocode pini1245) is a Pama-Nyungan language of the Wati group in Western Australia. It is spoken in the Northwest Goldfields, from around Barwidgee Station through Leinster and Wiluna to Sandstone. Wikipedia reports only about 6 fully fluent speakers. Some sources treat it as a dialect of the Western Desert Language, while Bowern classifies it as a distinct language. The closely related Birniridjara (``Pini'') is reportedly undocumented.

\section{Primary Educational Resources}
I found no primers, textbooks or school programmes for Tjupany.

\section{Linguistic and Reference Materials}
\begin{itemize}
\item Tjupan dialect/language (Wikipedia): \url{https://en.wikipedia.org/wiki/Tjupan_dialect}
\item Endangered Languages Project bibliography: \url{https://endangeredlanguages.com/lang/7069/bibliography}
\end{itemize}
The bibliography page may point to primary sources. I only saw its URL in the search results and did not read its contents, so I cannot name specific works.

\section{Religious and Community Resources}
Joshua Project lists a Tjupany people group of about 300 people. That is a population figure, not a count of speakers: \url{https://joshuaproject.net/people_groups/21574} and \url{https://joshuaproject.net/languages/tjp}. I found no community organization.

\section{Digital and Audio Resources}
None found.

\section{Curriculum Development Context}
Very little exists. With so few elderly speakers, urgent documentation should come first. Teaching could draw on materials for the related Western Desert Language, but this needs checking with community members and local language centres. I did not identify any in the search.
\end{document}